\section{Introduction}
\label{sec:introduction}

In March 2020, dollar funding markets seized. The cross-currency basis — the premium above the
covered-interest-parity rate that non-US institutions pay to borrow dollars in the FX swap market
— widened to levels not seen since the acute phase of the 2008 financial crisis.
The dislocations were not uniform across the global banking system. US commercial banks, whose
access to Federal Reserve emergency facilities is automatic and unconditional, were insulated.
Foreign banks whose central banks hold active swap lines with the Federal Reserve faced partial
disruption. Offshore dollar users — foreign financial institutions with no access to any official
backstop — bore the full force of the stress. This three-tier pattern is not a coincidence of March
2020; it repeats across every major dollar funding episode and reflects a structural feature of the
international monetary architecture: backstop coverage is hierarchical, and the hierarchy maps
directly to the distribution of funding stress. No existing theoretical model treats this
hierarchy as a formal primitive. This paper builds one.

The cross-currency basis literature and the international liquidity literature each capture part of
the picture, but each stops short of the mechanism that generates the hierarchy's economic
consequences. \citet{DuTepperVerdelhan2018} establish that CIP fails persistently for major
currencies since 2008 and attribute the deviation to balance-sheet constraints of financial
intermediaries, but their framework treats all foreign intermediaries symmetrically: the
differential between a bank whose central bank holds a swap line and one that does not plays
no formal role. \citet{BacchettaDavisVanWincoop2025} extend this to a general equilibrium
framework that jointly determines the spot rate and the CIP deviation through an endogenous
arbitrage capacity constraint $\chi_t$, yet this constraint does not derive from an explicit
backstop hierarchy — all CIP arbitrageurs face the same friction, and the swap-line coverage rate
$\lambda$ is absent from the model. \citet{CesaBianchiEgurenMartinFerrero2025} analyze Federal
Reserve swap lines as a mechanism that relaxes the funding constraint of a representative foreign
bank, but the treatment is single-segment and single-instrument: one class of foreign intermediary,
one official facility, with no feedback from fire-sale-driven Treasury price movements to dealer
wealth and hence to swap supply. At the conceptual level, \citet{MurauPapePforr2023} and
\citet{Mehrling2021} provide an institutional taxonomy of the international monetary hierarchy —
the Federal Reserve at the apex, primary dealers at an intermediate tier, foreign banks with swap
lines at a subordinate tier, and offshore dollar users without any backstop at the periphery — but
this architecture has no formal representation: there are no equilibrium conditions, no
optimization problems, and no mechanism linking backstop structure to the cross-currency basis.
The gap is precise. No model endogenizes backstop quality as a structural primitive that
simultaneously determines funding-market segmentation, the amplification of shocks, and the
directional propagation of stress across tiers.

This paper fills that gap. We build a static two-date model in which five classes of agents interact
in a dollar funding market structured by a three-tier backstop hierarchy. Segment 1 (US commercial
banks) holds a full backstop $B_1 = M_1$ and is inactive in equilibrium. Segment 2 (foreign banks
whose central banks hold swap lines) receives partial coverage $B_2 = \lambda M_2$, leaving a
residual funding need $R_2 = (1-\lambda)M_2$ to be met through market instruments. Segment 3
(offshore dollar users) has no official backstop, $B_3 = 0$, and must cover $R_3 = M_3 + \varepsilon$
entirely through market channels, where $\varepsilon \geq 0$ parameterizes funding stress. A single
global dealer intermediates the FX swap market and holds a Treasury inventory $\bar{X}_d$ whose
marked-to-market value evolves with the Treasury price: $W(P) = W_0 + P\bar{X}_d$. The dealer
operates under a binding Basel-type leverage constraint that limits total risk-weighted swap
exposure to a multiple of equity. This connects dealer swap capacity directly to Treasury prices
through the endogeneity of $W(P)$, closing the amplification loop. Outside investors absorb fire
sales from segments 2 and 3 along a downward-sloping demand curve. The equilibrium reduces to a
single scalar equation in the cross-currency basis $b^*$:
\begin{equation}
  b^* = \frac{R - \Gamma\tilde{W}}{\Delta},
  \label{eq:intro-eq}
\end{equation}
where $\Delta \equiv \Theta(P_0 - \Gamma\gamma\bar{X}_d)$ is a stability parameter, $\Gamma = \phi/q_S$
is the dealer capacity parameter, $\Theta$ is the aggregate fire-sale elasticity, $\gamma$ is the
outside-investor Treasury price elasticity, and $\tilde{W} = W_0 + \bar{X}_d$ is dealer wealth at
par. Existence and uniqueness follow from two conditions: $\Delta > 0$ (stability) and $R > \Gamma\tilde{W}$
(dollar scarcity). The architecture of backstop coverage, the dealer's leverage constraint, and the
fire-sale dynamics are all internal to the equilibrium.

This paper makes three formal contributions.

\textit{First}, we derive the amplification multiplier for the international dollar-funding market.
Proposition~2 shows that the general equilibrium basis decomposes as
$b^* = b^{\rm pe} \cdot \mathcal{M}$,
where $b^{\rm pe}$ is the partial-equilibrium basis absent dealer-wealth feedback and
$\mathcal{M} = 1/(1-\rho)$ with $\rho \equiv \Gamma\gamma\bar{X}_d/P_0$.
The ratio $\rho$ indexes the strength of the feedback loop running from fire sales to Treasury
prices, through dealer wealth, and back to swap supply. The multiplier $\mathcal{M}$ is the closed-form
international-dollar-funding analog of the liquidity spiral described in
\citet{BrunnermeierPedersen2009}: the same geometric series of feedback rounds summed in closed
form, but specialized to the cross-currency basis, FX swap supply, and the leverage constraint
of a global dealer. Proposition~3 characterizes the monotone dependence of $\mathcal{M}$ on each
structural primitive — dealer capacity $\Gamma$, outside-investor elasticity $\gamma$, dealer
inventory $\bar{X}_d$ — enabling direct mapping from institutional parameters to amplification
intensity.

\textit{Second}, we establish hierarchical contagion: stress in the uncovered segment propagates
upward through the backstopped segment to the dealer. Proposition~4 shows that a shock
$d\varepsilon > 0$ to segment 3 raises segment 2 fire sales by
$dx_2^*/d\varepsilon = \theta_2\mathcal{M}/(P_0\Theta) > 0$
despite $dR_2 = 0$ — segment 2 bears no direct funding shock yet responds through the price
mechanism because the equilibrium basis has risen. Dealer wealth contracts
($dW^*/d\varepsilon < 0$) and swap supply compresses ($dO^*/d\varepsilon < 0$). The contagion is
hierarchically ascending: it travels from the lowest tier upward. This result has no counterpart
in existing single-segment frameworks; it requires the three-tier structure as a primitive.

\textit{Third}, we derive the policy frontier for regulation and backstop provision. Proposition~6
establishes that tighter capital regulation (parameterized as an increase $\tau$ in the swap
risk-weight) raises the equilibrium basis under mild conditions, and that regulation and backstop
provision are imperfect substitutes along an explicit policy frontier. Define $\tau^{\rm safe}(\lambda, b^{\max})$
as the maximum $\tau$ consistent with $b^* \leq b^{\max}$. Then $\partial\tau^{\rm safe}/\partial\lambda > 0$:
a stronger swap-line commitment expands the set of regulatory tightening compatible with stress
below any given target. The trade-off $d\tau/d\lambda|_{b^*=\text{const}}$ is finite and well-defined,
but varies with $\rho$ in a way that diverges as $\rho \to 1$ — near the breakdown point, backstop
provision is essential and no amount of regulatory adjustment substitutes for it.

Our paper contributes to three strands of the literature. Among models of CIP deviation and
dollar funding stress, we advance beyond \citet{DuTepperVerdelhan2018} and
\citet{BacchettaDavisVanWincoop2025} by endogenizing the segmentation that these papers take as
parametric: the analog of Bacchetta et al.'s capacity constraint $\chi_t$ is derived here from
the three-tier backstop architecture rather than imposed. Among models of global intermediary
capacity, we extend \citet{GabaixMaggori2015} from a single homogeneous financier class to a
three-tier hierarchy with heterogeneous backstop coverage and an endogenous feedback from asset
prices to dealer wealth. Cross-segment contagion (Proposition~4) has no counterpart in their
framework. Among analyses of official liquidity facilities, we add a second channel to the
single-segment mechanism in \citet{CesaBianchiEgurenMartinFerrero2025}: Proposition~5 shows that
swap lines work not only through direct demand relief but also through a collateral-stabilization
channel that amplifies each dollar of support by the factor $\mathcal{M}$, and that this multiplier is
largest precisely during acute stress. The three-tier structure also formalizes the institutional
taxonomy of \citet{MurauPapePforr2023} and \citet{Mehrling2021}, converting their hierarchy from
a descriptive framework into an equilibrium model with closed-form predictions. The connection to
the inside/outside liquidity framework of \citet{HolmstromTirole1998} runs through Proposition~7:
the run-zone thresholds $(\underline\lambda, \bar\lambda)$ are expressed in terms of the structural
parameters of the dollar funding market, extending the \citet{DiamondDybvig1983} and
\citet{ChangVelasco2001} coordination-failure logic to the international backstop setting.

The paper proceeds as follows. \Cref{sec:model} develops the model environment, states the
optimization problems, and derives the equilibrium. \Cref{sec:results} presents the amplification
multiplier (Proposition~2) and its comparative statics (Proposition~3). \Cref{sec:contagion}
analyzes hierarchical contagion across segments (Proposition~4) and the double dividend of
swap lines (Proposition~5). \Cref{sec:policy} derives the regulation-backstop policy frontier
(Proposition~6). \Cref{sec:runs} extends the model to backstop uncertainty and self-fulfilling
runs (Proposition~7). \Cref{sec:conclusion} concludes.
